% =====================================================================
%  Research report — free LaTeX template by mitaa.dev
%  Compile with pdfLaTeX + Biber (Overleaf does this for you).
%  Locally:  latexmk -pdf main.tex
% =====================================================================
\documentclass[11pt, a4paper, oneside]{report}
% =====================================================================
%  Styling of the research report template — mitaa.dev
%  The content lives in main.tex. Change the font, colour and
%  margins here; a regular report needs none of it.
% =====================================================================

% --- Language and font ----------------------------------------------
\usepackage[T1]{fontenc}
\usepackage[english]{babel}
\usepackage{libertinus}          % Libertinus: serif, sans and maths
\usepackage{microtype}

% --- Table of contents (load before parskip) ------------------------------
\usepackage[titles]{tocloft}  % title set as an ordinary heading
\renewcommand{\cftchapfont}{\sffamily\bfseries}
\renewcommand{\cftchappagefont}{\sffamily\bfseries}
\setlength{\cftbeforechapskip}{0.7em}

% --- Page ------------------------------------------------------------
\usepackage[a4paper, margin=2.5cm]{geometry}
\usepackage{parskip}             % space between paragraphs, no indent
\linespread{1.08}

% --- Colour -----------------------------------------------------------
\usepackage{xcolor}
\definecolor{accent}{HTML}{1F5FA8}   % the one accent colour; change to taste
\definecolor{subtle}{HTML}{6B7280}

% --- Content ----------------------------------------------------------
\usepackage{amsmath}
\usepackage{graphicx}
\graphicspath{{figures/}}        % put your images in the figures/ folder
\usepackage{booktabs}
\usepackage{tabularx}
\usepackage{enumitem}
\setlist{itemsep=0.2em, topsep=0.3em}
\usepackage{caption}
\DeclareCaptionFont{accent}{\color{accent}}
\captionsetup{font=small, labelfont={bf,accent}, labelsep=period}
\usepackage{siunitx}             % numbers and units: \qty{9.81}{\metre\per\second\squared}

% --- References (APA 7) -------------------------------------------------
\usepackage{csquotes}
\usepackage[style=apa, backend=biber]{biblatex}
\addbibresource{references.bib}

% --- Headings ----------------------------------------------------------
\usepackage{titlesec}
\titleformat{\chapter}[display]
  {\sffamily\bfseries\color{accent}}
  {\small\MakeUppercase{\chaptertitlename}\ \thechapter}{0.5em}
  {\fontsize{24}{28}\selectfont\color{black}}
  [\vspace{0.4em}{\color{accent}\rule{2.5em}{2.5pt}}]
\titleformat{\section}{\sffamily\Large\bfseries\color{accent}}{\thesection}{0.75em}{}
\titleformat{\subsection}{\sffamily\large\bfseries}{\thesubsection}{0.75em}{}
\titlespacing*{\chapter}{0pt}{0pt}{2em}
\titlespacing*{\section}{0pt}{1.8em}{0.6em}
\titlespacing*{\subsection}{0pt}{1.3em}{0.3em}

% --- Report details ---------------------------------------------------
% Fill these in at the top of main.tex; they appear on the title page.
% \docfield{name} creates \docname{...} to fill in and \thedocname to print it.
\newcommand{\docfield}[1]{%
  \expandafter\newcommand\csname thedoc#1\endcsname{}%
  \expandafter\newcommand\csname doc#1\endcsname[1]{%
    \expandafter\renewcommand\csname thedoc#1\endcsname{##1}}}
\docfield{title}
\docfield{subtitle}
\docfield{authors}
\docfield{programme}
\docfield{course}
\docfield{supervisor}
\docfield{client}
\docfield{institution}
\docfield{version}
\docfield{date}
\docdate{\today}

% --- Header and footer ------------------------------------------------
\usepackage{fancyhdr}
\pagestyle{fancy}
\fancyhf{}
\fancyhead[L]{\small\sffamily\color{subtle}\thedoctitle}
\fancyhead[R]{\small\sffamily\color{subtle}\nouppercase{\leftmark}}
\fancyfoot[R]{\small\sffamily\thepage}
\renewcommand{\headrulewidth}{0pt}
\renewcommand{\chaptermark}[1]{\markboth{#1}{}}
\renewcommand{\sectionmark}[1]{}
\fancypagestyle{plain}{\fancyhf{}\fancyfoot[R]{\small\sffamily\thepage}}

% --- Title page -----------------------------------------------------------
% One line of details; empty fields are skipped.
\newcommand{\docmeta}[2]{%
  \ifdefempty{#2}{}{%
    \noindent\makebox[8em][l]{\color{subtle}#1}%
    \begin{tabular}[t]{@{}l@{}}#2\end{tabular}\par}}

\newcommand{\maketitlepage}{%
  \begin{titlepage}
    \raggedright\sffamily
    {\color{accent}\rule{3em}{3pt}}\par
    \vspace{2.5em}
    {\small\bfseries\color{accent}RESEARCH REPORT\par}
    \vspace{0.8em}
    {\fontsize{30}{35}\selectfont\bfseries\thedoctitle\par}
    \ifdefempty{\thedocsubtitle}{}{%
      \vspace{0.8em}{\Large\color{subtle}\thedocsubtitle\par}}
    \vfill
    {\setlength{\parskip}{0.35em}%
      \docmeta{Authors}{\thedocauthors}
      \docmeta{Programme}{\thedocprogramme}
      \docmeta{Course}{\thedoccourse}
      \docmeta{Supervisor}{\thedocsupervisor}
      \docmeta{Client}{\thedocclient}
      \docmeta{Institution}{\thedocinstitution}
      \docmeta{Version}{\thedocversion}
      \docmeta{Date}{\thedocdate}}
  \end{titlepage}}

% --- Abstract keywords ---------------------------------------------------
\newcommand{\keywords}[1]{\par\medskip\noindent{\sffamily\bfseries\color{accent}Keywords:}\ #1\par}

% --- Links (load last) -----------------------------------------
\usepackage[hidelinks]{hyperref}
\usepackage[nameinlink, noabbrev]{cleveref}


% --- Details: fill these in -------------------------------------------
\doctitle{The Title of Your Research}
\docsubtitle{A subtitle that sums up the research question}
\docauthors{First Last (1234567)}
\docprogramme{Degree programme}
\doccourse{Graduation research}
\docsupervisor{Supervisor's name}
\docclient{Client organisation}   % leave empty if there is none
\docinstitution{University or college}
\docversion{1.0}
\docdate{\today}

\begin{document}

\maketitlepage

% --- Front matter: Roman page numbers -------------------------------------
\pagenumbering{roman}

\chapter*{Preface}
A short, personal introduction: how this research came about and whom you
would like to thank. Sign with your name, place and date.

\bigskip
\noindent First Last\\
Place, \today

\chapter*{Abstract}
Summarise the research in 250 to 400 words: the background, the research
question, the method, the main results and the conclusion. Someone who
reads only this page should know what you studied and what came out of
it.

\keywords{keyword one, keyword two, keyword three}

\tableofcontents

% --- Main text: Arabic page numbers -----------------------------------------
\clearpage
\pagenumbering{arabic}

% =====================================================================
\chapter{Introduction}
\label{ch:introduction}
% =====================================================================

\section{Background}
Describe the situation that led to this research. Why does the subject
matter now, and to whom?

\section{Problem statement}
Name the problem as concretely as you can: what goes wrong or is
unknown, and what are the consequences?

\section{Objective}
The objective of this research is \ldots{} Phrase it so that, at the end,
you can check whether it was achieved.

\section{Research questions}
\label{sec:questions}
The main question of this research is:

\begin{quote}
  \itshape How \ldots{} ?
\end{quote}

To answer the main question, the following sub-questions were drawn up:
\begin{enumerate}[label=\arabic*.]
  \item What is already known about \ldots{} ?
  \item What does the current situation at \ldots{} look like?
  \item Which solutions are possible for \ldots{} ?
\end{enumerate}

\section{Outline}
\Cref{ch:theory} sets out the theoretical framework and \cref{ch:method}
the approach. \Cref{ch:results} answers the sub-questions, after which
\cref{ch:conclusion} answers the main question. \Cref{ch:discussion}
discusses the limitations and \cref{ch:recommendations} closes with
recommendations.

% =====================================================================
\chapter{Theoretical framework}
\label{ch:theory}
% =====================================================================
Describe the theories, models and earlier studies you build on. Cite your
sources in parentheses \parencite{creswell2018} or in the running text,
as with \textcite{booth2016}. Show how the literature sheds light on your
problem and where a gap remains.

\section{Key concepts}
Define the concepts you measure or study, so that the reader knows what
you mean by them.

\section{Conceptual model}
A conceptual model shows which factors you study and how they relate
(\cref{fig:model}).

\begin{figure}[htbp]
  \centering
  \includegraphics[width=0.6\linewidth]{example-image-b}
  \caption{Conceptual model of the research. Replace this example with
    your own figure in the \texttt{figures/} folder.}
  \label{fig:model}
\end{figure}

% =====================================================================
\chapter{Method}
\label{ch:method}
% =====================================================================

\section{Research design}
Explain which type of research you do (qualitative, quantitative or
mixed) and why it fits the research questions.

\section{Population and sample}
Describe whom or what you study, how you selected the participants or
objects, and how large the sample is.

\section{Data collection}
Describe the instruments: interviews, questionnaires, measurements or
document analysis. The questionnaire is in \cref{app:questionnaire}.

\section{Data analysis}
Explain how you processed the data. If you use a formula, number it, like
the mean in \cref{eq:mean}:
\begin{equation}
  \label{eq:mean}
  \bar{x} = \frac{1}{n} \sum_{i=1}^{n} x_i
\end{equation}

\section{Reliability and validity}
Describe what you did to make the results reliable and valid, for
example a pilot, triangulation or a second assessor.

% =====================================================================
\chapter{Results}
\label{ch:results}
% =====================================================================
Discuss the results per sub-question, in the same order as in
\cref{sec:questions}. Describe what you found, without explaining it yet.

\section{Sub-question 1}
\Cref{tab:descriptive} summarises the respondents' answers.

\begin{table}[htbp]
  \centering
  \caption{Descriptive statistics per group ($n = 120$).}
  \label{tab:descriptive}
  \begin{tabular}{l S[table-format=2.0] S[table-format=1.2] S[table-format=1.2]}
    \toprule
    Group & {$n$} & {Mean} & {Standard deviation} \\
    \midrule
    Group A & 40 & 3.82 & 0.71 \\
    Group B & 42 & 3.47 & 0.88 \\
    Group C & 38 & 4.05 & 0.64 \\
    \bottomrule
  \end{tabular}
\end{table}

\section{Sub-question 2}
Describe the results of the second sub-question. Put a quotation from an
interview in quotation marks: \enquote{That is simply how it always went}
(respondent 4).

\section{Sub-question 3}
Describe the results of the third sub-question.

% =====================================================================
\chapter{Conclusion}
\label{ch:conclusion}
% =====================================================================
First briefly answer each sub-question, then the main question. Use only
what is in \cref{ch:results}; new information does not belong here.

% =====================================================================
\chapter{Discussion}
\label{ch:discussion}
% =====================================================================
Interpret the results: what do they mean, and how do they relate to the
literature in \cref{ch:theory}? Discuss the limitations of the research
and what they mean for the conclusion. Close with suggestions for further
research.

% =====================================================================
\chapter{Recommendations}
\label{ch:recommendations}
% =====================================================================
Give concrete, actionable recommendations for the client. For each one,
say what should happen, by whom and why.

\printbibliography[title={References}, heading=bibintoc]

% --- Appendices -------------------------------------------------------------
\appendix

\chapter{Questionnaire}
\label{app:questionnaire}
Include the full questionnaire or interview guide here.

\chapter{Raw data}
Put the raw data or extended tables here that would make the main text too
long.

\end{document}
